\documentclass[11pt]{article}

\usepackage[a4paper,margin=1in]{geometry}
\usepackage{amsmath,amssymb}
\usepackage{hyperref}

\hypersetup{
    colorlinks=true,
    linkcolor=blue,
    citecolor=blue,
    urlcolor=blue
}

\title{A Simple Example of Brenier's Theorem}
\author{}
\date{}

\begin{document}

\maketitle

Consider two probability measures on $\mathbb R$:
\[
\mu=\mathrm{Unif}[0,1],
\qquad
\nu=\mathrm{Unif}[1,2].
\]
The optimal transport map for the quadratic cost
\[
c(x,y)=|x-y|^2
\]
is
\[
T(x)=x+1.
\]
Indeed, $T$ pushes $\mu$ forward to $\nu$:
\[
T_\#\mu=\nu,
\]
because translating a uniform random variable on $[0,1]$ by $1$ gives a
uniform random variable on $[1,2]$.

\section*{Brenier Form}

Brenier's theorem says that, under its hypotheses, the optimal map has the form
\[
T=\nabla\varphi
\]
for a convex function $\varphi$.

In one dimension, the gradient is just the derivative. Here
\[
T(x)=x+1
\]
is the derivative of
\[
\varphi(x)=\frac12 x^2+x.
\]
Since
\[
\varphi''(x)=1>0,
\]
the function $\varphi$ is convex. Therefore
\[
T(x)=\varphi'(x)=x+1
\]
is exactly a Brenier map.

\section*{Interpretation}

This example says that the optimal way to move the uniform distribution on
$[0,1]$ to the uniform distribution on $[1,2]$ is to translate every particle
one unit to the right.

\end{document}
